\chapter{Numerical Linear Algebra and Performance}
\label{chap:numerics}

\Cref{chap:pumphb} specified the algorithm. This chapter treats the linear
algebra that makes it tractable: Schur reduction, factorisation backends, matrix
ordering, and the measured consequences of each choice. Every performance number
quoted here was measured on the machine described in \cref{sec:hardware}, not
estimated.

\section{Schur reduction}
\label{sec:schur}

\subsection{The structure to exploit}

A TWPA has thousands of nodes but only a handful matter as observables: the
ports. Moreover the vast majority of nodes are \emph{linear} --- they carry no
nonlinear branch --- so their equations are exactly solvable given the states of
their neighbours.

Partition the node set into \emph{retained} nodes $\mathcal{R}$ (ports, and any
node touching a nonlinear branch) and \emph{eliminated} nodes $\mathcal{E}$. The
dynamic block partitions conformally:
\begin{equation}
  \Dmat_k =
  \begin{bmatrix}
    \Dmat_{nn,k} & \Dmat_{ne,k}\\
    \Dmat_{en,k} & \Dmat_{ee,k}
  \end{bmatrix},
\end{equation}
and the linear rows give
\begin{equation}
  X_{e,k} = \Dmat_{ee,k}^{-1}\left(\Svec_{e,k} - \Dmat_{en,k}X_{n,k}\right).
  \label{eq:back-substitution}
\end{equation}
Substituting into the retained rows leaves a reduced system in $X_{n,k}$ alone,
governed by the Schur complement
\begin{equation}
  \boxed{\;
  \mathcal{S}_k = \Dmat_{nn,k} - \Dmat_{ne,k}\,\Dmat_{ee,k}^{-1}\,\Dmat_{en,k}.
  \;}
  \label{eq:schur-complement}
\end{equation}

Everywhere $\Dmat_k$ appeared in \cref{chap:pumphb}, $\mathcal{S}_k$ appears in
the reduced problem, and the algorithm is otherwise unchanged.

\subsection{Why this is cheap here}

Two properties make \eqref{eq:schur-complement} inexpensive:

\begin{enumerate}
\item $\Dmat_{ee,k}$ is \textbf{constant in $\Xvec$}. The eliminated nodes are
      linear by construction, so the factorisation
      $\mat{P}_r\Dmat_{ee,k}\mat{P}_c = \mat{L}_{e,k}\mat{U}_{e,k}$ is performed
      \emph{once per frequency and harmonic} and reused beneath the power,
      continuation, Newton, and GMRES loops. This is the first of the two LU
      factorisations of \cref{rem:two-lus}.
\item $\Dmat_{en,k}$ is \textbf{extremely sparse in its columns}. Only retained
      columns adjacent to an eliminated node contribute; for a ladder that is a
      tiny interface set.
\end{enumerate}

\begin{implnote}[column-block assembly with fill dropping]
\code{assemble\_schur\_complements} solves only the active columns:
\begin{lstlisting}
active = np.unique(Den.nonzero()[1])
W_active = part.dee_factors[h].solve(Den[:, active].toarray())
\end{lstlisting}
No dense inverse of $\Dmat_{ee}$ is ever formed. Entries below
\code{drop\_tol} $=10^{-12}$ times the per-block maximum magnitude are dropped,
because the numerical fill produced by the dense solve is at round-off level and
would otherwise destroy the banded sparsity that makes the reduced problem cheap.

A \code{matrix\_free} mode is also available, applying
$\mathcal{S}_k\vect{v} = \Dmat_{nn,k}\vect{v} - \Dmat_{ne,k}(\Dmat_{ee,k}^{-1}
(\Dmat_{en,k}\vect{v}))$ through the prefactored block without assembling
$\mathcal{S}_k$; the assembled form is the default because it can also serve as
a preconditioner.
\end{implnote}

\begin{caveat}[the retained set must be well posed]
If $\Dmat_{ee,k}$ is singular the elimination is invalid --- the retained set is
too small to determine the eliminated nodes. The builder raises with exactly
that message rather than returning a silently wrong complement.
\end{caveat}

\subsection{Consequence for warm starting}

Because the engine's state $\Xvec$ is Schur-reduced, it has the
\emph{retained-port shape}, which is constant across frequencies. Chained warm
starts, residual ranking, and power substeps all operate on the same
representation, with no disk round-trip and no shape conversion. A pump solution
loaded from disk is full-shape and must be promoted; that is a separate code
path.

\section{Factorisation backends}

\subsection{PARDISO versus SuperLU}

Two sparse direct solvers are available for the preconditioner factorisation.

\begin{measured}[JTWPA, $S=10$]
\begin{center}
\begin{tabular}{@{}l r r@{}}
\toprule
Backend & Refactor time & Peak RSS \\
\midrule
PARDISO (\code{pypardiso} 0.4.7) & \SI{0.92}{\second} & \SI{3.01}{\giga\byte} \\
SuperLU (\code{scipy.sparse.linalg.splu}) & \SI{3.94}{\second} & \SI{2.81}{\giga\byte} \\
\bottomrule
\end{tabular}
\end{center}
PARDISO is $4.3\times$ faster for $\SI{7}{\percent}$ more peak memory. Since
peak memory sets the worker count and both give the same count here, PARDISO
wins by $\approx 3.8\times$ at equal concurrency.
\end{measured}

The gap comes from symbolic reuse. This problem refactors at every Newton step
with an unchanging sparsity pattern, and PARDISO exposes a phase-22
(numeric-only) call that reuses the analysis. \code{spla.splu} has no
symbolic-reuse API, so SuperLU repeats the ordering every time.

\begin{implnote}[fail loudly during timed runs]
The fallback to SuperLU is automatic and silent. Setting
\code{TWPA\_REQUIRE\_PARDISO=1} makes it a hard error, so a timed campaign
cannot silently produce $4\times$-slow numbers because an optional dependency
failed to import.
\end{implnote}

\subsection{Thread count}

\begin{implnote}[PARDISO pinned to one thread]
\code{TWPA\_PARDISO\_THREADS} defaults to $1$. The reason is robustness: MKL
intermittently fails reordering with error $-3$ at its default thread count in
this mixed OpenBLAS + MKL environment.

The pin costs single-run latency, and the scaling is pathological on the AMD
part used here:
\begin{center}
\begin{tabular}{@{}r r r@{}}
\toprule
Threads & Factor time ($S=6$) & Relative \\
\midrule
1 & \SI{340}{\milli\second} & $1.00\times$ \\
2 & \SI{228}{\milli\second} & $1.49\times$ \\
6 & \SI{1374}{\milli\second} & $0.25\times$ \\
\bottomrule
\end{tabular}
\end{center}
At six threads MKL PARDISO is four times \emph{slower} than serial. Raising the
count only helps when running a single worker; with several workers, extra
workers beat extra threads --- measured $3\times1$ thread gives
\SI{5.00}{cycles\per\second} against $3\times2$ threads at
\SI{3.66}{cycles\per\second}.
\end{implnote}

\subsection{Banded reordering}
\label{sec:banded}

A TWPA is a one-dimensional chain, and that structure can be exposed by
reordering.

The coupled Jacobian of \cref{sec:real-coupled} is naturally packed
\emph{tone-major}: index $= (\text{tone})\times n + (\text{node})$. In that
ordering the matrix spans the full dimension, because every tone couples to
every other tone at every node. Reordering \emph{node-major},
\begin{equation}
  \text{index} = (\text{node})\times 2H + (\text{super-block}),
\end{equation}
collapses the matrix onto a band roughly three tone-blocks wide, because the
\emph{circuit} couples only nearest-neighbour nodes. The factors then fit a
LAPACK general band rather than a general sparse LU.

\begin{remark}[this is a property of the circuit, not of an ordering search]
The bandwidth reduction follows from the chain topology and is not something a
generic bandwidth-reducing permutation discovers better. Reverse Cuthill--McKee
does \emph{worse}: maximum bandwidth 147 against 89 on JTWPA at $S=2$. The
bandwidth is nonetheless \emph{measured} from the assembled pattern, never
assumed.
\end{remark}

\begin{measured}[JTWPA $S=10$, three power points, single process]
\begin{center}
\begin{tabular}{@{}l r r r r@{}}
\toprule
Backend & Peak RSS & Wall & Workers in \SI{7}{\giga\byte} & Net throughput \\
\midrule
\code{pardiso} & \SI{2.51}{\giga\byte} & \SI{147.7}{\second} & 2 & $1.55\times$ \\
\code{banded}  & \SI{1.84}{\giga\byte} & \SI{173.7}{\second} & 3 & $\mathbf{1.95\times}$ \\
\bottomrule
\end{tabular}
\end{center}
Converged gain agrees to $\SI{4.4e-10}{\dB}$
($27.541036124719$ against $27.541036124280$), as it must --- this changes only
the preconditioner, not the equation.
\end{measured}

The trade is $\SI{18}{\percent}$ more wall time for $\SI{27}{\percent}$ less
memory. That is worth it only when the smaller footprint buys another worker.
For a single-frequency run it is a pure $1.18\times$ loss, so the default
remains \code{pardiso} and the frequency-resolved campaigns pass \code{banded}
explicitly.

\section{Assembly cost}

\subsection{The scatter map}

Assembling the real-coupled matrix \eqref{eq:real-superblock} at every Newton
step is itself expensive if done naively. The pattern is fixed by the basis, so
only the values change.

\begin{implnote}[from sparse scatter matrix to index arrays]
The original design held the assembly as a sparse matrix $\mat{W}$ of shape
$(\nnz(\Mmat),\,2 n_\ell \nnz_k)$, making assembly a single sparse
matrix--vector product. That was elegant but wasteful: every stored value is
$\pm1$ and the source indices are contiguous per-$\ell$ ranges, so $\mat{W}$
carried no information beyond target indices --- and cost
$\approx\SI{630}{\mega\byte}$ at $S=10$ to hold \SI{94}{\mega\byte} of indices.

It was replaced by four \code{int32} target arrays per (mode, mode) block
(\code{\_index\_contributions}). Peak RSS fell from \SI{3.04}{\giga\byte} to
\SI{2.51}{\giga\byte} for $+\SI{87}{\milli\second}$ of assembly time ---
scattered writes replacing a sequential-write product.
\end{implnote}

\begin{measured}[cached assembly plus symbolic reuse, per Newton step]
\begin{center}
\begin{tabular}{@{}l r r r@{}}
\toprule
Path & Before & After & Speed-up \\
\midrule
Full backend (fixture) & \SI{6.148}{\second} & \SI{0.93}{\second} & $6.6\times$ \\
Schur backend (circuit dir) & \SI{5.004}{\second} & \SI{1.00}{\second} & $5.0\times$ \\
\bottomrule
\end{tabular}
\end{center}
The assembled matrix is identical to the reference assembly, verified to
$\num{1.0e-11}$ relative.
\end{measured}

\subsection{Memory scaling}

\begin{proposition}[memory scales as the square of the tone count]
The coupled Jacobian is block-dense in tone index: every retained tone couples
to every other through $\Khat_{k\pm q}$. Memory therefore scales as
\begin{equation}
  \left(n_{\mathrm{pump\;modes}} + 2S + 1\right)^{2},
\end{equation}
not as the packed dimension.
\end{proposition}

\begin{measured}[same circuit, $n=2048$]
\begin{center}
\begin{tabular}{@{}l r r r@{}}
\toprule
& Tones $H$ & $\nnz(\Mmat)$ & Dimension \\
\midrule
Gain-map pump & 10 & \num{2.46e6} & $\times 1.0$ \\
Multitone $S=10$ & 31 & \num{23.6e6} & $\times 3.1$ \\
\midrule
Ratio & $3.1\times$ & $\mathbf{9.6\times}$ & --- \\
\bottomrule
\end{tabular}
\end{center}
$9.6 \approx (31/10)^{2}$: exactly the $H^{2}$ law, while the dimension grew
only $3.1\times$.
\end{measured}

Per-worker peak memory follows: $\approx\SI{2.80}{\giga\byte}$ at $S=10$,
$\approx\SI{1.6}{\giga\byte}$ at $S=6$, $\approx\SI{0.9}{\giga\byte}$ at $S=2$.

\begin{implnote}[worker count is computed, not guessed]
\code{multitone.resources.fast\_coupled\_footprint} estimates the per-worker
footprint from the basis size and caps \code{--signal-workers} against
\emph{both} the declared resource budget and actual free memory. The previous
hardcoded \SI{3.0}{\giga\byte} per worker was independent of basis size,
underestimated $S=10$ by $\SI{38}{\percent}$, and caused an out-of-memory
failure at two workers on a \SI{15.3}{\giga\byte} machine.
\end{implnote}

\section{Preconditioner reuse}
\label{sec:precond-reuse}

Modified-Newton preconditioning reuses one factorisation across $N$ consecutive
Newton steps. Since the update is always taken against the \emph{true} Jacobian
(only the preconditioner is stale), the converged solution cannot change --- only
the iteration count.

\begin{measured}[JPA compression fixture]
\begin{center}
\begin{tabular}{@{}r r r r@{}}
\toprule
$N$ & Factorisations saved & Extra GMRES iterations & Wall time \\
\midrule
1 & --- & --- & \SI{3313}{\milli\second} \\
2 & 26 of 59 & $+841$ ($243\to1084$) & \SI{3574}{\milli\second} \\
3 & 29 of 59 & $+986$ & \SI{3545}{\milli\second} \\
\bottomrule
\end{tabular}
\end{center}
Gain identical to all reported digits in every case.
\end{measured}

\begin{remark}[why reuse cannot pay here]
The standard argument for modified Newton assumes a preconditioner that is
cheap relative to a Krylov iteration, so amortising it over several steps is a
win. That regime does not exist in this problem. The exact preconditioner
converges GMRES in about three iterations, and one GMRES iteration --- a
Jacobian-vector product plus a triangular solve, \SI{303}{\milli\second} at
$S=6$ --- costs about as much as one factorisation (\SI{346}{\milli\second}).
There is nothing to amortise, and a stale exact factor is much worse than a
fresh one.

The default is therefore $N=1$: refactor every step. The pump solve is
additionally pinned to $N=1$ so that its iterate path is bit-for-bit unchanged
by any experiment with this setting.
\end{remark}

\section{Worker scaling}

\begin{measured}[JTWPA $S=2$, six-core part]
\begin{center}
\begin{tabular}{@{}r r@{}}
\toprule
Workers & Relative throughput \\
\midrule
1 & $1.00\times$ \\
2 & $1.55\times$ \\
3 & $2.30\times$ \\
4 & $2.00\times$ \\
6 & $2.41\times$ \\
\bottomrule
\end{tabular}
\end{center}
Run-to-run spread $\pm\SI{15}{\percent}$.
\end{measured}

Scaling is markedly sub-linear and plateaus around three workers: the problem is
memory-bandwidth-bound, not compute-bound. Sparse triangular solves and
scattered assembly writes are both bandwidth-limited operations, so adding cores
past the point where the memory controller saturates buys little.

\section{Map traversal and recovery}

A gain map is a two-dimensional grid over pump power and pump frequency, and the
order in which cells are visited determines how good the warm starts are.

\paragraph{Traversal orders.} \code{column} (the legacy per-frequency pass),
\code{backbone}, \code{nearest}, \code{serpentine}, and \code{floodfill}. The
non-column orders share one in-process store across both axes, so they force
single-process execution.

\paragraph{Recovery ladders.} On a failed cell: \code{reseed} (fresh cold
solve), \code{alt\_parent} (warm start from a different neighbour),
\code{bridge} (physical-parameter continuation from a solved parent along
$(P,f)$), \code{ladder}.

\paragraph{Power substep.} The most effective recovery is adaptive continuation
along the map's own power axis. On a failed warm cell, the solver walks up in
adaptive \si{\dBm} micro-steps from the last converged state --- geometric in
current, growing $\times1.5$ on success and halving on failure --- warm-starting
each.

\begin{measured}[power substep on a real column]
At $f_p = \SI{8.099}{\giga\hertz}$, the coarse \SI{0.30}{\dB} power grid failed
at $\SI{-29.36}{\dBm}$ with 3--4 cells converged. Substepping recovered to
$\approx\SI{-25.75}{\dBm}$: 16 cells converged, $+\SI{3.9}{\dB}$ of additional
range, passing through multi-lobe gain ripple that the coarse grid had
overshot, before reaching a genuine boundary.
\end{measured}

A step-independent stall --- the step shrinking below its minimum with no
progress --- is \emph{recorded} as a numerical or fold boundary rather than
retried indefinitely.

\begin{caveat}[terminology]
Throughout, a failure to converge is called a \emph{convergence failure}, not a
``fold'', unless the physics has been verified. The two are genuinely different:
a fold is a property of the solution branch, a convergence failure is a property
of the solver. Conflating them was a real source of confusion in earlier work.

Evidence that at least one boundary is genuine: a diagnostic that forces the
gain solve on the \emph{last Newton iterate regardless of convergence} shows the
gain \emph{collapsing} from \SI{15.7}{\dB} to \SI{1.5}{\dB} at the first
non-converged point. The non-converged waveform does not sustain the gain, which
is what one expects near a real transition rather than a mere solver miss.
\end{caveat}

\section{Hardware and reproducibility}
\label{sec:hardware}

All timings quoted in this chapter were measured on a six-core AMD part with
approximately \SI{7}{\giga\byte} of usable RAM for solver processes, under
Windows, with \code{pypardiso} 0.4.7, SciPy 1.18.0 and NumPy, PARDISO pinned to
one thread. Memory figures are peak resident set size of a single worker
process.

Two consequences worth stating: the worker-count ceilings quoted are specific to
this memory budget and scale directly with available RAM; and the PARDISO thread
pathology is specific to this MKL-on-AMD combination and should be re-measured
on different hardware rather than assumed.
